% ============================================================
\subsection{SQL Patterns (PostgreSQL)}
% ============================================================

Reference for recurring query shapes. Each entry states when to reach for
it and the mechanism that makes it work.

\begin{keybox}[Logical processing order]
A query is evaluated in this order, regardless of written order:
\[
  \texttt{FROM/JOIN} \to \texttt{WHERE} \to \texttt{GROUP BY} \to
  \texttt{HAVING} \to \texttt{SELECT} \to \texttt{DISTINCT} \to
  \texttt{ORDER BY} \to \texttt{LIMIT}.
\]
Two consequences drive most of the patterns below. \texttt{WHERE} runs
before aggregation, so it cannot see aggregates or \texttt{SELECT} aliases.
Window functions are computed at the \texttt{SELECT} step, so they also
cannot appear in \texttt{WHERE} and must be wrapped in a subquery or CTE
to filter on.
\end{keybox}

\subsubsection{Filtering and Joins}

\paragraph{WHERE vs HAVING.} \texttt{WHERE} filters individual rows before
grouping. \texttt{HAVING} filters whole groups after aggregation and is the
only place an aggregate predicate is legal. A filter on a plain
(non-aggregated) column belongs in \texttt{WHERE} so rows are dropped before
the grouping work happens. Reserve \texttt{HAVING} for conditions like
\texttt{COUNT(*) > 5}.

\vspace{1.5ex}

\begin{lstlisting}[style=sql, caption={WHERE prunes rows, HAVING prunes groups}]
SELECT department, COUNT(*) AS headcount
FROM employees
WHERE active = true          -- row filter, before grouping
GROUP BY department
HAVING COUNT(*) > 5;         -- group filter, on an aggregate
\end{lstlisting}

\vspace{-2ex}

\paragraph{Joins.}
\begin{itemize}[label=$\circ$]
  \item \textbf{INNER:} rows with a match on both sides.
  \item \textbf{LEFT / RIGHT:} keep all rows from one side, \texttt{NULL}-fill the other.
  \item \textbf{FULL OUTER:} keep unmatched rows from both sides.
  \item \textbf{CROSS:} Cartesian product, every left row against every right row.
  \item \textbf{SELF:} a table joined to itself under an alias (employee to manager, row to prior row).
  \item \textbf{Semi-join} (\texttt{EXISTS}, \texttt{IN}): keep left rows with at least one match. No row duplication, no right-side columns.
  \item \textbf{Anti-join} (\texttt{NOT EXISTS}, or \texttt{LEFT JOIN ... WHERE right.key IS NULL}): keep left rows with no match.
\end{itemize}

\subsubsection{Window Functions}

Every window function shares the same shape:
\texttt{f() OVER (PARTITION BY ... ORDER BY ...)}. \texttt{PARTITION BY} splits
the rows into independent groups, \texttt{ORDER BY} sets the order within each
group, and the function is evaluated per row without collapsing the result the
way \texttt{GROUP BY} does. All of them are computed at the \texttt{SELECT}
step, so to filter on their output you wrap them in a subquery or CTE.

\paragraph{Rank-then-filter with ROW\_NUMBER().} Use this for top-$N$ per
group when you need to keep the full row. \texttt{GROUP BY} would collapse the
other columns away. Number rows inside each partition, then filter on that
number in an outer query, since the window function cannot live in
\texttt{WHERE}.

Choose the ranker by tie policy: \texttt{ROW\_NUMBER} assigns distinct numbers
and breaks ties arbitrarily, \texttt{RANK} repeats a rank and leaves gaps,
\texttt{DENSE\_RANK} repeats with no gaps.

\begin{lstlisting}[style=sql, caption={Highest-paid employee per department}]
SELECT department, name, salary
FROM (
    SELECT department, name, salary,
           ROW_NUMBER() OVER (PARTITION BY department
                              ORDER BY salary DESC) AS rn
    FROM employees
) ranked
WHERE rn = 1;
\end{lstlisting}

\paragraph{Offset access with LAG() and LEAD().} These read a value from
another row relative to the current one, without a self-join.
\texttt{LAG(col, k, default)} returns the value \texttt{k} rows earlier in the
ordered partition, \texttt{LEAD} looks \texttt{k} rows ahead. Both default to
an offset of \texttt{k = 1} and return \texttt{NULL} past the partition edge
unless a fill \texttt{default} is supplied. The \texttt{ORDER BY} inside
\texttt{OVER} defines what ``previous'' means and is mandatory.
\texttt{PARTITION BY} resets the sequence at each group boundary, so a lag
never bleeds from one partition into the next.

The dominant use is a row-to-row comparison: a delta, a period-over-period
change, or a return. Dividing the current value by its lag gives a growth
factor.

\begin{lstlisting}[style=sql, caption={Daily returns per ticker with LAG}]
SELECT ticker, trade_date, close,
       close / LAG(close) OVER (PARTITION BY ticker
                                ORDER BY trade_date) - 1 AS daily_return
FROM prices;
\end{lstlisting}

\vspace{-1.5ex}

The first row of each partition has no predecessor, so its \texttt{LAG} is
\texttt{NULL} and the derived return is \texttt{NULL}. That is correct here,
since there is no prior price to compare against. \texttt{LEAD} is the same
tool aimed forward, used below for gap detection, where a jump between an id
and the next present id marks a break in a run.

\paragraph{Running and rolling aggregates.} An ordinary aggregate gains a
window when you attach \texttt{OVER (... ORDER BY ...)}. With an
\texttt{ORDER BY} and no explicit frame, the frame defaults to everything from
the partition start up to the current row, giving a running (cumulative)
total. Add an explicit \texttt{ROWS BETWEEN} frame to get a fixed moving
window instead.

\begin{lstlisting}[style=sql, caption={20-day moving average and running volume}]
SELECT ticker, trade_date, close,
       AVG(close) OVER (PARTITION BY ticker ORDER BY trade_date
                        ROWS BETWEEN 19 PRECEDING AND CURRENT ROW) AS ma_20,
       SUM(volume) OVER (PARTITION BY ticker ORDER BY trade_date)  AS cum_volume
FROM prices;
\end{lstlisting}

\texttt{ma\_20} is the current row plus the 19 before it, a 20-row moving
average, and \texttt{cum\_volume} uses the default frame for a running total.
These are the SQL forms of a pandas \texttt{.rolling(20).mean()} (fixed
\texttt{ROWS} frame) and \texttt{.expanding().sum()} (running total from the
start).

One subtlety: the default frame is \texttt{RANGE UNBOUNDED PRECEDING}, which
includes every row tied on the \texttt{ORDER BY} key, not just rows physically
up to the current one. When ties on the order key exist and you want a
row-precise window, state \texttt{ROWS BETWEEN} explicitly rather than relying
on the default.

\paragraph{Gaps and islands.} An \emph{island} is a maximal run of consecutive
values. The trick: subtract a dense row number from the value. Across a
consecutive run both sides increase by one per step, so the difference stays
constant, and it jumps exactly at a gap. Group by the difference to collapse
each run.

\begin{lstlisting}[style=sql, caption={Islands of consecutive ids}]
WITH numbered AS (
    -- DISTINCT first so duplicate ids do not desync the row number
    SELECT DISTINCT log_id,
           ROW_NUMBER() OVER (ORDER BY log_id) AS rn
    FROM logs
)
SELECT MIN(log_id) AS island_start,
       MAX(log_id) AS island_end,
       COUNT(*)    AS length
FROM numbered
GROUP BY log_id - rn
ORDER BY island_start;
\end{lstlisting}

The \texttt{DISTINCT} is load-bearing. If \texttt{log\_id} repeats, \texttt{rn}
still advances on the duplicate while the value does not, so
\texttt{log\_id - rn} shifts and splits an island that should stay whole. For
consecutive \emph{dates} rather than integers, offset by an interval:
\texttt{log\_date - rn * INTERVAL `1 day'} is constant within a run.

\begin{lstlisting}[style=sql, caption={Gaps via LEAD}]
SELECT log_id + 1  AS missing_from,
       next_id - 1 AS missing_to
FROM (
    SELECT log_id,
           LEAD(log_id) OVER (ORDER BY log_id) AS next_id
    FROM logs
) t
WHERE next_id - log_id > 1;   -- values missing between the two present ids
\end{lstlisting}

\newpage

\subsubsection{Aggregation}

\paragraph{Conditional aggregation (pivot).} Fold a categorical column into one
row of counts by wrapping a \texttt{CASE} in an aggregate. Rows the
\texttt{CASE} does not match contribute 0 to \texttt{SUM}. In Postgres the
\texttt{FILTER} clause does the same thing and reads more cleanly.

\begin{lstlisting}[style=sql, caption={Pivot risk categories to columns}]
SELECT
    inspection_type,
    COUNT(*) FILTER (WHERE risk_category IS NULL)           AS no_risk,
    COUNT(*) FILTER (WHERE risk_category = 'Low Risk')      AS low_risk,
    COUNT(*) FILTER (WHERE risk_category = 'Moderate Risk') AS moderate_risk,
    COUNT(*) FILTER (WHERE risk_category = 'High Risk')     AS high_risk,
    COUNT(*)                                                AS total
FROM sf_restaurant_health_violations
GROUP BY inspection_type
ORDER BY total DESC;
\end{lstlisting}

The \texttt{SUM(CASE WHEN ... THEN 1 ELSE 0 END)} form is equivalent and
portable across engines. \texttt{COUNT} ignores \texttt{NULL}, so
\texttt{COUNT(CASE WHEN cond THEN 1 END)} (no \texttt{ELSE}) also counts only
the matches.

\subsubsection{NULL Handling}

\paragraph{COALESCE and NULLIF.} \texttt{COALESCE(a, b, c, ...)} returns the
first non-\texttt{NULL} argument, evaluated left to right, or \texttt{NULL} if
all are \texttt{NULL}. The common use is a fallback that keeps a \texttt{NULL}
from poisoning arithmetic: \texttt{salary + commission} is \texttt{NULL}
whenever \texttt{commission} is, so write \texttt{salary + COALESCE(commission, 0)}.
It also supplies a default down a priority list, as in
\texttt{COALESCE(nickname, first\_name, 'Unknown')}. It is ANSI standard,
unlike the vendor-specific \texttt{IFNULL} (MySQL) and \texttt{ISNULL}
(SQL Server), and it accepts any number of arguments.

\texttt{NULLIF(a, b)} is the mirror image: it returns \texttt{NULL} when
\texttt{a = b}, otherwise \texttt{a}. Its standard use guards division,
\texttt{x / NULLIF(y, 0)} yields \texttt{NULL} instead of raising a
divide-by-zero when \texttt{y} is 0.

\begin{lstlisting}[style=sql, caption={NULLIF guards division, COALESCE supplies defaults}]
SELECT
    order_id,
    revenue / NULLIF(quantity, 0)  AS unit_price,    -- NULL, not an error, when quantity = 0
    COALESCE(discount, 0)          AS discount,       -- NULL discount treated as 0
    COALESCE(nickname, first_name) AS display_name
FROM orders;
\end{lstlisting}

A subtlety with aggregates: \texttt{SUM}, \texttt{AVG}, and
\texttt{COUNT(col)} skip \texttt{NULL}s. So \texttt{AVG(commission)} averages
only the non-\texttt{NULL} rows, while \texttt{AVG(COALESCE(commission, 0))}
counts the \texttt{NULL}s as zeros and enlarges the denominator. Decide which
you mean before wrapping the column.

\newpage

\paragraph{Gotchas.}
\begin{itemize}
  \item \textbf{\texttt{NOT IN} with NULLs:} if the subquery returns any
  \texttt{NULL}, \texttt{x NOT IN (...)} is never true and the query returns
  zero rows. \texttt{x NOT IN (1, NULL)} expands to
  \texttt{NOT (x=1 OR x=NULL)}, and the \texttt{NULL} comparison poisons the
  result to unknown. Use \texttt{NOT EXISTS}, which is null-safe.
  \item \textbf{Integer division truncates:} \texttt{5 / 2} yields \texttt{2}
  because the fractional part is discarded. Cast one operand:
  \texttt{5::numeric / 2} or \texttt{5.0 / 2}. Bites hardest in ratio and
  percentage columns.
  \item \textbf{\texttt{COUNT(col)} vs \texttt{COUNT(*)}:} the former skips
  \texttt{NULL}s in \texttt{col}, the latter counts every row. A common source
  of off-by-$N$ errors.
\end{itemize}